% Part: lambda-calculus
% Chapter: syntax
% Section: terms

\documentclass[../../../include/open-logic-section]{subfiles}

\begin{document}

\olfileid{lam}{syn}{trm}
\olsection{पदे}

लॅम्डा कलनातील पदे $\Obj{v_0}$, $\Obj{v_1}$, \dots या
अमर्याद चरांपासून, ~“$\lambd$” या चिन्हापासून आणि कंसांपासून
विगमनाधारित रीतीने घडवली जातात. चरे दर्शवण्यासाठी आपण
$x$, $y$, $z$, \dots{} आणि पदे दर्शवण्यासाठी
$M$, $N$, $P$, \dots{} वापरू.

\begin{defn}[पदे] \ollabel{defn:term}
लॅम्डा कलनातील \emph{पदांचा} संच पुढीलप्रमाणे
विगमनाधारित रीतीने परिभाषित करतात:
\begin{enumerate}
  \item \ollabel{defn:term-var} $x$ हे चर असेल, तर $x$ हे
    पद आहे.
  \item \ollabel{defn:term-abs} $x$ हे चर आणि $M$ हे
    पद असेल, तर $(\lambd[x][M])$ हे पद आहे.
  \item \ollabel{defn:term-app} $M$ आणि $N$ ही दोन्ही
    पदे असतील, तर $(MN)$ हे पद आहे.
\end{enumerate}
\end{defn}

पद $(\lambd[x][M])$ हे \olref{defn:term-abs} नुसार
घडवले असेल, तर ते \emph{अमूर्तीकरणाचे} फलित आहे
असे म्हणतो; आणि $\lambd[x]$ मधील $x$ ला
\emph{!!{parameter}} म्हणतात. \olref{defn:term-app}
नुसार घडवलेले पद $(MN)$ हे \emph{उपयोजनाचे}
फलित असते.

वर परिभाषित पदांमध्ये सर्व कंस स्पष्टपणे लिहिलेले
असतात. उदाहरणार्थ,
$(\lambd[x][((\lambd[x][x])(\lambd[x][(xx)]))])$
या पदावरून दिसते की हे खूपच अवजड होऊ शकते. कंस
टाळण्यासाठी आपण काही रूढी मांडू. मात्र अधिकृत
व्याख्येमुळे \olref{defn:term} नुसार पद कसे
घडले हे ठरवणे सोपे जाते. उदाहरणार्थ,
$(\lambd[x][((\lambd[x][x])(\lambd[x][(xx)]))])$
हे पद घडवण्याची शेवटची पायरी अमूर्तीकरणच असली
पाहिजे; त्यात !!{parameter} हे~$x$ आहे.
ते $((\lambd[x][x])(\lambd[x][(xx)]))$ या पदावर
अमूर्तीकरण करून मिळते; हे पद दोन पदांच्या
उपयोजनातून घडले आहे. त्या दोन पदांपैकी प्रत्येक
अमूर्तीकरणातून घडले आहे; आणि पुढे असेच चालते.

\begin{prob}
$(\lambd[g][(\lambd[x][(g (x x))])
  (\lambd[x][(g (x x))])])$ हे पद कसे घडले ते वर्णन करा.
\end{prob}

\end{document}
